\section{Introduction}
\label{sec:introduction}

The sieve of Eratosthenes repeatedly removes multiples of known primes
\cite{hardywright1979}. A standard implementation stores a bounded array
and crosses out composites. The Sieve Sequence studied here exposes a
different mathematical view: after a finite set of prime filters has been
installed, the survivors form an infinite periodic sequence. A finite
cycle of adjacent gaps therefore represents the whole stage.

This representation belongs to the broad family of cyclic or wheel-based
sieve descriptions. Pritchard's survey places wheel sieves among
systematic families of prime-number sieves \cite{pritchard1987}. The
present contribution is not the wheel idea itself. It is the explicit
decomposition of the stage and transition into contracts that can be
checked by Stainless, a verifier for Scala programs
\cite{stainless2024verification}, against a first-principles arithmetic
and list library.

The proof is organized into these property groups:

\begin{itemize}
  \item stage semantics and complete increasing enumeration --- Section~3;
  \item canonical period and finite gap-cycle reconstruction --- Section~4;
  \item repeated-cycle invariance, exact filtering, and copy-or-merge
    dynamics --- Section~5;
  \item next-head primality and next-stage agreement --- Section~6;
  \item conditional assumptions and open composition problems ---
    Section~7.
\end{itemize}

The mathematical object is separate from its proof namespaces.
\texttt{SpecSieveSequence} is the data model and linear semantic
specification. Independent property objects establish the period,
survivor count, transition, next-stage assembly, and head-primality
theorems.
